\PassOptionsToPackage{backend=bibtex}{biblatex}
\documentclass[9pt,a4paper,twocolumn,twoside]{rho-class/rho}
\usepackage[english]{babel}
\usepackage{float}

%----------------------------------------------------------
% Title
%----------------------------------------------------------

\doctype{Proyecto 2. Seminario de Heurísticas de Optimización Combinatoria.}
\title{Implementación y experimentación con MBO para JSSP}

%----------------------------------------------------------
% Authors, affiliations and dates
%----------------------------------------------------------

\author[{\textasteriskcentered},a]{Wendy Sánchez Cortés}
\affil[a]{Facultad de Ciencias, Universidad Nacional Autónoma de México}
\dates{Fecha de entrega: \today}

\journal{Reporte del Proyecto 2}
\theday{\today}
\vol{1}
\no{1}

%----------------------------------------------------------
% Abstract and keywords
%----------------------------------------------------------

\begin{abstract}
    El presente documento describe la implementación y el proceso de experimentación
    para resolver una instancia particular del Problema de Programación de Talleres de
    Trabajo (JSSP) con 40 trabajos y 40 máquinas, donde cada trabajo consta de 40 subtareas.
    Para ello se utilizó la heurística Migrating Birds Optimization, propuesta por Ekrem Duman,
    Mitat Uysal y Ali Fuat Alkaya.
\end{abstract}

\begin{document}

\maketitle

\section{Introducción}


Para este segundo proyecto del Seminario de Heurísticas de Optimización Combinatoria se decidió implementar,
de forma voluntaria, la heurística Migrating Birds Optimization (MBO). El criterio de elección, su simplicidad.
Inspirados en la formación en V de las aves migratorias, Duman, Uysal y Alkaya publicaron en 2012
"Migrating Birds Optimization: A new metaheuristic approach and its performance on quadratic assignment problem".\\

El sistema conserva la idea del proyecto anterior: una interfaz abstracta \texttt{Solution} que desacopla la lógica
de la heurística del problema concreto, por lo que MBO no sabe nada de máquinas ni de trabajos. La implementación está
en C++, evalúa las soluciones con el algoritmo de Kahn, ejecuta los experimentos en paralelo con múltiples hilos y
genera reportes en formato CSV.

\section{Background: Recocido Simulado}

\section{Diseño del Sistema}
\subsection{Representación del problema}
\subsection{Arquitectura del sistema}
\section{Experimentación y resultados}

\section{Dificultades y Limitaciones}

\section{Conclusiones}

\end{document}